\chapter{Operational Trust as the Limiting Resource of Digital-Twin
Deployments}\label{ch:trust}

\section{Research question}\label{sec:trust-rq}

This chapter addresses RQ2: \emph{what limits the operational authority
a digital twin is granted -- model accuracy, or something else?} The
question matters because deployment decisions are usually argued in
accuracy terms, while the corpus examined here suggests the binding
constraint lies elsewhere.

\section{The argument}\label{sec:trust-argument}

Three strands of evidence converge on the same answer: the limiting
resource is \emph{operator trust}, and trust responds to the twin's
honesty about its own state rather than to its accuracy in the mean.

First, the synchronisation literature separates the two directly. In
comparative trials, explicit staleness marking -- returning every value
with its age, so a twin that has stopped tracking reality says so --
roughly halved operator mistrust incidents without any change to the
underlying model {[}1{]}. Accuracy was held constant; honesty moved;
trust followed the honesty.

Second, the longest observational record in the corpus shows authority
following audited agreement rather than claimed accuracy. The packaging
line twin of {[}2{]} gained scheduling authority station by station, and
each extension was argued from a published override log -- supervisors
observed agreeing with the twin in over ninety-five percent of cases --
not from validation statistics. The two failures that record contains
were both \emph{silent} divergences between plant and model, and the
study's operational recommendation is to treat any such divergence as a
severity-one defect {[}2{]}.

Third, the same asymmetry appears at the analytics layer. Comparing
anomaly detectors over twin state streams, the learned model with the
best mean precision was also the one whose failure was invisible: it
absorbed a slow drift into its own retraining baseline and missed it
entirely, while an incorruptible statistical baseline flagged it
{[}3{]}. The recommended layered configuration buys accuracy where it is
cheap but anchors trust to the component whose errors are legible.

\section{Synthesis}\label{sec:trust-synthesis}

Across all three strands, the mechanism that converts a working twin
into a trusted one is the same: make its uncertainty visible at the
point of use, and let humans audit the record on their own schedule. The
vocabulary of {[}4{]} sharpens the claim: what distinguishes a twin from
a shadow is a closed actuation loop, and closing that loop is precisely
the step organisations refuse until the audit trail exists. Authority,
in other words, is granted against the override log, not against the
validation report.

\section{Limitations}\label{sec:trust-limitations}

The corpus is small and synthetic, its trust measures are behavioural
proxies (mistrust incidents, override rates) rather than instruments,
and no study in it manipulates honesty and accuracy independently. The
claim this chapter defends is therefore directional: where the two have
been separated at all, honesty dominated.

\phantomsection\label{refs}
\begin{CSLReferences}{0}{0}
\bibitem[\citeproctext]{ref-sample_dt_sync_2023}
\CSLLeftMargin{{[}1{]} }%
\CSLRightInline{C. Chen and D. Devi, {``State synchronisation strategies
for operational digital twins,''} \emph{Synthetic Sample Papers}, vol.
1, pp. 6--11, 2023.}

\bibitem[\citeproctext]{ref-sample_dt_factory_2022}
\CSLLeftMargin{{[}2{]} }%
\CSLRightInline{E. Eriksen, {``A digital twin on the factory floor: An
eighteen-month case study,''} \emph{Synthetic Sample Papers}, vol. 1,
pp. 12--17, 2022.}

\bibitem[\citeproctext]{ref-sample_ml_anomaly_2023}
\CSLLeftMargin{{[}3{]} }%
\CSLRightInline{F. Farah and G. Gupta, {``Lightweight anomaly detection
over digital-twin state streams,''} \emph{Synthetic Sample Papers}, vol.
1, pp. 18--23, 2023.}

\bibitem[\citeproctext]{ref-sample_dt_overview_2024}
\CSLLeftMargin{{[}4{]} }%
\CSLRightInline{A. Author and B. Builder, {``Digital twins: Definitions,
distinctions, and a short taxonomy,''} \emph{Synthetic Sample Papers},
vol. 1, pp. 1--5, 2024.}

\end{CSLReferences}
